
En este apéndice se listan los repositorios de código desarrollado como parte del trabajo y las dependencias de terceros utilizadas (también sus repositorios). Salvo que se indique lo contrario, se usa la rama \texttt{main} con la versión más reciente de cada repositorio. La única excepción es \ac{rtxi}, del que se usó un \textit{tag} específico.

\section{Código desarrollado\label{SEC:CODIGODESARROLLADO}}

\subsection{Código fuente del proyecto\label{SEC:PARAMETRIZACIONREPO}}

En este repositorio se encuentra el código fuente desarrollado en este trabajo, incluyendo la implementación \textit{offline} y el módulo \textit{online} para \ac{rtxi} del algoritmo de parametrización mediante \ac{bo}, y el módulo de neurona de \ac{hr} modificada para \ac{rtxi} corregido a partir del repositorio de \ac{rthybrid} para \ac{rtxi} (Apéndice~\ref{SEC:RTXIPLUGINSGNB}).

\url{https://github.com/marcogomezhernandez2004/tfg-code}


\subsection{Datos del proyecto\label{SEC:DATOSREPO}}

En este repositorio se encuentran todos los datos, gráficas y códigos involucrados (códigos de simulación y generación de gráficas y datos) del proyecto. Incluye contenido adicional al presentado en este trabajo.

\url{https://github.com/marcogomezhernandez2004/tfg-data}


\section{Repositorios de terceros (dependencias)\label{SEC:REPOSTERCEROS}}

\subsection{Neun\label{SEC:NEUNREPO}}

Librería C++ para la simulación de sistemas dinámicos y redes neuronales, desarrollada por el \ac{gnb} de la \ac{uam}. Se usa en la versión \textit{offline} para implementar la neurona postsináptica y la sinapsis. Se descarga automáticamente al compilar.

\url{https://github.com/GNB-UAM/Neun}


\subsection{\acs{rthybrid}\label{SEC:RTHYBRIDREPO}}

Librería de modelos de neuronas y sinapsis para circuitos electrofisiológicos biohíbridos en \ac{rt}, desarrollada por el \ac{gnb} de la \ac{uam}. El algoritmo de escalado de la señal presináptica se basa en su implementación.

\url{https://github.com/GNB-UAM/RTHybrid}


\subsection{\acs{rthybrid} para \acs{rtxi}\label{SEC:RTXIPLUGINSGNB}}

Conjunto de módulos de \ac{rtxi} desarrollados por el \ac{gnb} de la \ac{uam} que proporcionan la funcionalidad de \ac{rthybrid}, incluyendo módulos de neuronas modelo, análisis de \textit{bursts} y escalado de amplitud. Se instalan como módulos independientes en \ac{rtxi}.

\url{https://github.com/GNB-UAM/rthybrid-for-rtxi}


\subsection{\texttt{mimic-signal}\label{SEC:MIMICSIGNALREPO}}

Módulo estándar de \ac{rtxi} que copia una señal de entrada y la emite con ganancia y \textit{offset} configurables. Se usa para aplicar el escalado de amplitud a la señal presináptica, debiendo introducir los términos de escalado manualmente. Se instala como módulo independiente en \ac{rtxi}.

\url{https://github.com/RTXI/mimic-signal}


\subsection{\acs{rtxi}\label{SEC:RTXIREPO}}

Plataforma de experimentación en \ac{hrt} de código abierto sobre la que se ejecuta la versión \textit{online}, en este caso con un \textit{kernel} Linux parcheado con Xenomai. Se usó la versión con el \textit{tag} \texttt{v2.3} (\textit{commit} \texttt{b8b5ffe}) por compatibilidad con la instalada en el laboratorio del \ac{gnb} de la \ac{uam}, debiendo hacer uso de la plantilla para módulos de \ac{rtxi} antigua (Apéndice~\ref{SEC:PLUGINTEMPLATE}).

\url{https://github.com/RTXI/rtxi}


\subsection{Plantilla de módulo de \texttt{rtxi} (\texttt{plugin-template})\label{SEC:PLUGINTEMPLATE}}

Plantilla base oficial para la creación de módulos de \ac{rtxi}. En la implementación \textit{online} se utilizó la versión antigua de esta plantilla (ficheros \texttt{plugin-template}) para asegurar la compatibilidad con la versión 2.3 de \ac{rtxi} (Apéndice~\ref{SEC:RTXIREPO}).

\url{https://github.com/RTXI/plugin-template}


\subsection{KFR\label{SUBSEC:KFRREPO}}

Biblioteca C++ de procesamiento digital de señal. Se usa para el filtrado Butterworth mediante \texttt{filtfilt} y las operaciones vectoriales sobre las señales (mediante el tipo \texttt{univector}). Se descarga automáticamente al compilar en la versión \textit{offline}; en la versión \textit{online} debe estar preinstalada en el sistema, requiriendo compilarse previamente con la bandera \texttt{-DCMAKE\_POSITION\_INDEPENDENT\_CODE=ON} y ajustar manualmente la arquitectura en el \texttt{Makefile} (cambiando \texttt{-lkfr\_dsp\_avx2} por la correspondiente al sistema). Esto se explica en el \texttt{README.md} de la versión \textit{online} (Apéndice~\ref{SEC:DATOSREPO}).

\url{https://github.com/kfrlib/kfr}


\subsection{Limbo\label{SUBSEC:LIMBOREPO}}

Biblioteca C++ de \ac{bo} de optimización Bayesiana y modelado con procesos gaussianos, utilizada para ejecutar la optimización de los parámetros sinápticos. Se descarga automáticamente al compilar en la versión \textit{offline}; incluida como \textit{git submodule} en la versión \textit{online} (carpeta \texttt{extern}).

\url{https://github.com/resibots/limbo}


\subsection{nlohmann/json\label{SUBSEC:JSONREPO}}

Biblioteca C++ de cabecera única para JSON, utilizada para el guardado del historial de la \ac{bo}. Se descarga automáticamente al compilar en la versión \textit{offline}; vendorizada como cabecera en el repositorio de la versión \textit{online} (carpeta \texttt{include}).

\url{https://github.com/nlohmann/json}


\subsection{yaml-cpp\label{SUBSEC:YAMLCPPREPO}}

Biblioteca C++ para la lectura y escritura de ficheros YAML, utilizada en la versión \textit{offline} para la configuración y la escritura de los resultados. Se descarga automáticamente al compilar.

\url{https://github.com/jbeder/yaml-cpp}


\subsection{NLopt\label{SUBSEC:NLOPTREPO}}

Biblioteca de optimización no lineal, usada por Limbo (Apéndice~\ref{SUBSEC:LIMBOREPO}) para la maximización de la función de adquisición mediante Subplex. Dependencia del sistema preinstalada en ambas versiones.

\url{https://github.com/stevengj/nlopt}

